\documentclass[11pt,a4paper]{article}
\usepackage{luatexja}
\usepackage{amsmath}
\usepackage{amssymb}
\usepackage{amsfonts}
\usepackage{bm}
\usepackage{geometry}
\geometry{left=25mm,right=25mm,top=25mm,bottom=25mm}
\usepackage{hyperref}

\title{双四元数二重時空におけるねじれ不整合としての量子重力}

\author{https://github.com/hypernumbernet}
\date{\today}

\begin{document}

\maketitle

\begin{abstract}
一般相対性理論は重力を単一の時空連続体の曲率として成功裏に記述するが、量子力学との調和は依然として捉えがたく、特異点、連続体仮説、および曲率の微視的起源の欠如に悩まされている。本論文は、伴う論文で導入された二重時空理論（DST）の量子セクターを展開する。そこでは時空は共有される多様体ではなく、すべての質量粒子に内在的に付随し、16実次元の双四元数代数 $\mathbb{H} \oplus \mathbb{H} \cong \mathrm{Cl}(3,1)$ に符号化された、コンパクト化されたミンコフスキー時空の対である。

許容構成上、ねじれスカラーは $|J|\le 3\pi^2/8$、同値に $|J_{\mathrm{norm}}|\le 1$ を満たす。平方 $+1$ の生成子から冪等射影 $P_L$、$P_R$、$P_{\mathrm{spin}}$ が得られ、$(1\pm i)/2$（$i^2=-1$）および合成 $P_{\mathrm{spin}}P_R$ は冪等でない。通常と双対の生成子は同軸では可換であり、軸外では必ずしも可換でない。巡回生成子は $IJ=K$ を満たし、$\mathbb{C}^2$ 上で $-i\sigma_a$ として作用する。すべての双対ラピディティで双対ローターは $\mathrm{SU}(2)$ のロドリゲス行列であり、ユニタリで行列式は $1$ である。$\mathrm{Cl}(3,1)$ の次数1生成子は $\{\gamma^\mu,\gamma^\nu\}=2\eta^{\mu\nu}$（$\eta=\mathrm{diag}(-1,+1,+1,+1)$）を満たし、$\gamma^0$ は $P_{\mathrm{spin}}$ に用いる双曲生成子 $j$ ではない。キリング自己対は不定であり、ボルン確率ではない~\cite{dst-d4l}。$\mathrm{Cl}(3,1)$ と $\mathrm{Cl}(9,1)$ は同型でなく、双対スピノル空間は十次元時空のマヨラナ--ワイル加群ではない。TEGR とアインシュタイン--ヒルベルト作用の変分的等価は、ローター–テトラッド対応を固定したうえで文献から採用する。運動量空間のディラック方程式は \(\mathrm{Cl}(3,1)\) におけるミンコフスキー二次形式の一次因子であり、質量殻上では双対スピノルが静止系の対 \((u,iu)\) を回復し、その対の縦ブーストはブースト軸に沿うカイラル因子の後の同一の殻上構成である。書かれた反対符号作用は不整合を自由にし共通ラピディティを源とし、保存するヤコビ積分は不定である。質量をねじれとする読み、ヒッグス機構、測定、経路積分は幾何的仮説として残す。
\end{abstract}

\section{はじめに：量子重力への代数的アプローチ}

一般相対性理論と量子力学の調和は、理論物理学における最も深遠な未解決問題の一つであり続けている。数十年にわたる努力にもかかわらず、両理論の経験的成功と整合しつつ概念的非両立性を解消する合意的枠組みは現れていない。ループ量子重力は幾何を離散化するが背景多様体と曲率を基本的なものとして保持する。弦理論は直接的な観測的支持なしに余剰次元と超対称性を導入する。漸近安全および因果力学的三角形分割は連続体極限と時空そのものの起源に苦闘する。既存のすべてのアプローチは共通の基礎的仮定を共有する。時空連続体仮説である。時空は物質から独立した滑らかで微分可能な多様体として存在し、重力はそのエネルギー--運動量を源とする曲率から生じると想定される。

この仮定は巨視的スケールでは極めて成功しているが、量子レベルでは深刻な困難をもたらす。曲がった背景上の紫外発散、ブラックホール情報パラドックス、宇宙論的特異点、そして——最も根本的には——なぜ、いかにして物質が時空を曲げるべきかを説明する物理機構が全く存在しないことである。マッハの原理は満たされないままである。時空は物質の分布によって決定される関係的構造ではなく、自律的実体として振る舞う。

本論文では異なるアプローチを展開する。基礎となる代数は $\mathrm{Cl}(3,1)$ と同型な16実次元双四元数代数であり、以下の恒等式に背景多様体やリーマン曲率は不要である。物理描像はローター対 $R_\text{usual}$ と $R_\text{dual}$ である。双対セクターの運動学は $\mathbb{C}^2$ 上の表現であり、重力現象論は $J=\frac12\sum_a(\alpha_a^2-\beta_a^2)$ で測るねじれ不整合である。$J$ とテレパラレルスカラーの同一視、およびアインシュタイン--ヒルベルト等価性は、ローター–テトラッド対応を固定したうえでTEGR文献から採用する。

すべての双対ラピディティで双対ローターは $\mathrm{SU}(2)$ のロドリゲス行列である。ユニタリで行列式 $1$、相対双対ローターも群に留まる。これは第\ref{sec:quantumstates}節に記す $\mathbb{C}^2$ 上の行列表現であり、ローター多様体をリーマン三次元球面と同一視することではない。符号 $(-,+,+,+)$ のディラック行列は $\mathrm{Cl}(3,1)$ の次数1生成子であり、$\{\gamma^\mu,\gamma^\nu\}=2\eta^{\mu\nu}$ を満たし、$P_{\mathrm{spin}}$ に用いる双曲バイベクトルではない。ボルン則はキリング自己対ではない。その対は符号 $(3,3)$ をもち、純楕円では負になる~\cite{dst-d4l}。位置および運動量演算子、時空上のディラック偏微分方程式、質量としてのねじれはここでは導かない。運動量空間のディラック方程式は、第\ref{sec:dirac}節に記すとおり、ミンコフスキー二次形式のクリフォード平方から従う。

連続体仮説を完全に棄却することにより、本理論は古典一般相対性理論の特異点を自動的に解消する。双対時空のコンパクト性はローター角を離散化する。許容構成上では代数的上限は $|J|\le 3\pi^2/8$（同値に $|J_{\rm norm}|\le 1$）であり、構成が許容のままである限り、$r=0$ またはビッグバンにおいて現れるであろう非有界曲率を禁止する。事象地平も同様に存在しない。光線は共鳴トンネルを介して有限のねじれ層を横断する。一般相対性理論からの強場におけるずれは原理上観測可能となり、本枠組みは双対ローター位相のコヒーレント制御を通じた重力工学への具体的経路を開く。

本論文の残りの構成は以下のとおりである。第\ref{sec:spinor}節では平方 $+1$ の射影を構成し、合成 $P_{\mathrm{spin}}P_R$ が冪等でないことを示す。第\ref{sec:quantumstates}節では双対生成子を $-i\sigma_a$ と同一視し、すべてのラピディティにおける双対ローターを $\mathrm{SU}(2)$ 行列と同一視する。キリング対はボルン則ではない。第\ref{sec:dirac}節では $\mathrm{Cl}(3,1)$ のクリフォード関係を述べ、$\gamma^0$ と双曲生成子 $j$ を区別し、ミンコフスキー二次形式を運動量空間のディラック方程式へ因数分解し、双対スピノルから質量殻上の解を組み立て、静止系の対の縦ブーストをその構成と同一視する。第\ref{sec:mass}--\ref{sec:dynamics}節では質量の値、ヒッグス機構、測定、経路積分を幾何的仮説として論じ、書かれた粒子作用が不整合を自由にし共通ラピディティを源とすることを示す。第\ref{sec:predictions}節は許容上限に縛られた予想である。第\ref{sec:comparison}節では内部双対と十次元時空超対称を書き分ける。第\ref{sec:experiments}節は近未来の実験である。第\ref{sec:status}節は量子セクターの代数的描像をまとめる。第\ref{sec:conclusions}節は要約である。

量子重力が $\mathrm{Cl}(3,1)$ から完結して従うとは主張しない。本論文の結果は16次元代数の恒等式であり、連続体の動力学は古典的近似のままである。

\section{双四元数スピノルと極小左イデアル}
\label{sec:spinor}

双四元数代数 $\mathbb{H} \oplus \mathbb{H}$ は、符号 $(-,+,+,+)$ のミンコフスキー時空のクリフォード代数 $\mathrm{Cl}(3,1)$ と標準的に同型であり、ディラックスピノルの自然かつ最小の実現を与える。スピノルを外部の表現空間として導入する従来のアプローチとは異なり、ここではスピノルは代数そのものの元として直接創発する。

既約表現への射影を与える原始冪等元は、平方が $+1$ の生成子から構成する必要がある。$g^2=+1$ ならば
\[
\frac{1+g}{2}
\]
は冪等である。軸 $0$ の双曲生成子を $j:=e_0 e_1$ と書く（$j^2=+1$）。この $j$ は時間ベクトル $\gamma^0=e_0$ ではない。符号 $(-,+,+,+)$ では $(\gamma^0)^2=-1$ である（第\ref{sec:dirac}節）。作業用の射影は
\[
P_{\mathrm{spin}} = \frac{1+j}{2}=\frac{1+e_0 e_1}{2}
\]
および相補的な空間カイラル対
\[
P_L = \frac{1-e_1}{2},\qquad P_R = \frac{1+e_1}{2}
\]
（$e_1^2=+1$）であり、$P_L^2=P_L$、$P_R^2=P_R$、$P_L+P_R=1$ を満たす。

積
\[
P = P_{\mathrm{spin}}\, P_R
\]
は冪等でない。$e_0 e_1$ と $e_1$ は反可換であり、因子は可換でなく $(P_{\mathrm{spin}}P_R)^2\neq P_{\mathrm{spin}}P_R$ である。可換な平方 $+1$ 対は $e_1$ と軸外ブースト $e_0 e_2$ であり、その積射影は冪等で、生成する左イデアルの実次元は $8$ である。スピノル加群としての既約性は未解決とする。

これに対し、$\mathrm{Cl}(3,1)$ の擬スカラー $i$（$i^2=-1$）だけから作った素朴な式
\[
\frac{1\pm i}{2}
\]
は冪等でない：$\bigl((1-i)/2\bigr)^2=-i/2$。平方が $+1$ になる合成子（例：$i\gamma_1\gamma_2$）は構成部品として許されるが、裸の擬スカラーは許されない。

候補となるスピノル空間の一般元は明示的に
\[
\psi = \bigl( a_0 + a_1 i + b_1 kI + b_2 kJ + b_3 kK + c_1 jI + c_2 jJ + c_3 jK \bigr) P,
\]
の形をとり、$a_0,a_1,b_i,c_i \in \mathbb{R}$、$P$ は平方 $+1$ の射影である（$P_{\mathrm{spin}}P_R$ ではない）。擬スカラー $i$ を量子力学の虚数単位と同一視するのは辞書であり、担体がディラック加群であることの証明ではない。双対写像 $X \mapsto X i$ はなお時間配向を反転する。代数的射影子そのものは $P_L$、$P_R$ であり、$(1\pm i)/2$ ではない。

ローターの代数元への作用は左乗算である。同軸の通常および双対生成子は可換であり、軸外の対は必ずしも可換でない。とくに $\Gamma_a$ はすべての通常セクター生成子 $i\Gamma_b$ と可換ではない。したがって $R_{\mathrm{total}}=R_{\mathrm{usual}}R_{\mathrm{dual}}$ のカイラル因子化は任意の軸に対する恒等式ではない。

双一次形式 $\bar{\psi}\psi\propto\operatorname{Tr}(\Omega_{\mathrm{biv}})$ および標準模型のフェルミオンセクター全体が16次元代数から従うという主張はここでは確立しない。代数から従うのは上の射影の算法である。

\section{量子状態としての双対時空回転}
\label{sec:quantumstates}

双対ローター
\[
R_\text{dual} = \exp\left( \sum_{a=1}^3 \frac{\phi_a}{2} \Gamma_a \right),
\]
（ここで $\Gamma_1 = I$、$\Gamma_2 = J$、$\Gamma_3 = K$ は四元数関係 $\Gamma_a^2 = -1$ および $IJ=K$ を満たす）は $\mathbb{C}^2$ 上で $-i\sigma_a$ により表現される。別の辞書 $I\leftrightarrow i\sigma$ は $IJ=K$ と両立しない。

すべての双対ラピディティ $\beta$ について指数はロドリゲス公式
\[
U(\beta)
=\cos(\lVert\beta\rVert/2)\,I
+\sin(\lVert\beta\rVert/2)\,n\cdot(-i\sigma)
\]
であり、$n=\beta/\lVert\beta\rVert$（$\beta\neq 0$）、$U(0)=I$ である。その行列はユニタリで行列式 $1$ である。相対双対ローターは $\mathrm{SU}(2)$ に留まる。逆は反対のラピディティなので、相対双対ローターは群の積である。その指標は二つの双対軸の整列を記録し、非可換な残余はトレース零の傾きである。計算軸まわりの $2\pi$ 双対回転は $-I$ である。双対時空のコンパクトなトポロジー（$\phi_a \bmod 4\pi$）はこの $\mathbb{C}^2$ 上の行列群と両立する。多様体の独立な微分同相 $S^3\cong\mathrm{SU}(2)$ はここでは証明しない。

位置および運動量演算子、正準交換関係、不確定性限界 $\Delta\phi\cdot\Delta\theta\ge 1/2$ は幾何的仮説である。

確率解釈はキリング自己対ではない。同一視
\[
|\psi|^2 \propto \tfrac12 B(\Omega_{\mathrm{biv}},\Omega_{\mathrm{biv}})
\]
は成り立たない。キリング対は符号 $(3,3)$ をもち、純楕円配置で負になる。したがって自己重なりは確率ではない。ボルン則はヒルベルト層に属し、$J$ の帰結ではない~\cite{dst-d4l}。

測定、もつれ、収縮としてのねじれ同期も未解決とする。本論文における双対セクターの運動学は、巡回生成子の $\mathbb{C}^2$ 表現を意味する。

\section{内在的超対称性としての双対時空}
\label{sec:susy}

宇宙定数問題は、真空の破局として広く知られ、理論物理学における最も深刻な微調整危機の一つをなす。量子場の理論は、真空エネルギー密度が紫外カットオフまでのあらゆる場のモードの零点揺らぎから生じると予測する。カットオフをプランクスケールにとると、これは \(M_\text{Pl}^4 \approx 10^{74}\,\text{GeV}^4\) 程度のエネルギー密度を与え、SI単位ではおよそ \(10^{113}\,\text{J/m}^3\) に換算される。しかし、現在の加速膨張から推論される観測値は \(\rho_\text{vac}^\text{obs} \approx 5.4 \times 10^{-10}\,\text{J/m}^3\) にすぎず、宇宙定数 \(\Lambda \sim 1.1 \times 10^{-52}\,\text{m}^{-2}\) に対応する。したがって理論予測と測定値の比は
\[
\frac{\rho_\text{vac}^\text{theory}}{\rho_\text{vac}^\text{obs}} \sim 10^{120}.
\]
に達する。既知の繰り込み手続きは、裸のパラメータと量子補正の間の \(10^{120}\) 分の1の精度での並外れた相殺なしにこの食い違いを除去できない。問題は技術的ではなく概念的である。固定された背景上の量子場の理論の真空に重力を整合的に結合させることは、観測される宇宙論を破壊せずにはできない。

二重時空理論は、超対称性がすでにすべての質量粒子が運ぶ16実次元双四元数代数の内部に\emph{内在的かつ不可分に}実現されていることを明らかにすることにより、この危機を解消する。新しい場や粒子は導入されず、外部の対称性も課されたり破られたりしない。

基礎となるローレンツ候補 \(\mathfrak{so}(3,1)\)（十二次元の和 \(\mathfrak{so}(3,1)\oplus\mathfrak{so}(3,1)\) ではない）の六つの生成子は、二次形式の符号によって区別される二つのセクターに分割される。すなわち生成子の平方が \(+1\) である双曲セクターと、生成子の平方が \(-1\) である巡回セクターである。同一軸の生成子は可換であり、軸外の交換子は必ずしも消滅しない。完全ローター \(R_\text{total}\) は両セクターに同時に作用する。双対写像 \(X \mapsto Xi\) は時間配向を反転しつつ符号を除いて二つのセクターを交換する。この交換は代数そのものの内部自己同型をなす。それは粒子を異なる種へと写すのではなく、同一の16次元対象の永続的に結合した二つの側面を単に交換する。したがって、任意のボソン励起の超対称パートナーはすでに双対セクター成分として存在する。発見すべき外部の超対称パートナーは存在しない。

スカラー \(J\) は不整合バイベクトル \(\Omega_\text{biv}\) から構成されるパラメータ空間二次形式 \(J=\frac12\sum_a(\alpha_a^2-\beta_a^2)\) である。キリング形式が双曲セクターで \(+8\)、巡回セクターで \(-8\) という反対の値をとるため、ボソン的寄与とフェルミオン的寄与は反対符号で入る。二つのセクターが整列しているとき、これらの寄与は正確に相殺する。\(J\) の残余真空値は双対セクター角の量子揺らぎによってのみ生成される。巡回生成子の平方は \(-1\) であり、双対時空を楕円的にする。そのローター多様体はコンパクトな三次元球面 \(S^3 \cong \mathrm{SU}(2)\) である。したがって角は周期的であり離散的に標本化される。大きな整数 \(N\) に対して \(\theta_a \in (2\pi/N)\mathbb{Z}\) である。コンパクト性は \(N^{-2}\) 程度の代数的上限を\emph{強制しない}。許容構成上では \(|J|\le 3\pi^2/8\)、あるいは \(J_{\rm norm}=\frac8{3\pi^2}J\) として \(|J_{\rm norm}|\le 1\) である。離散量 \(\varepsilon_N=16/(3N^2)\) は非零正規化高さの\emph{下限}であり、上限ではない。静止でないすべての許容格子点は \(|J_{\rm norm}|\ge\varepsilon_N\) を満たす。真空抑制が起こるとすれば、それはセクター間の動力学的相殺から来なければならず、\(|J|\le J_{\max}\sim N^{-2}\) と書くことによるものではない。

双対ローターがすべてのラピディティで \(\mathrm{SU}(2)\) のロドリゲス行列であるため、双対セクターはスピン \(1/2\) の運動学を担い得る。標準模型フェルミオンと双対位相の完全な同一視は主張しない。双対写像 \(X\mapsto Xi\) は代数の内部自己同型である。それは十次元弦の時空超対称ではない：\(\mathrm{Cl}(3,1)\not\cong\mathrm{Cl}(9,1)\)、双対スピノル空間は \(\mathrm{Cl}(9,1)\) のマヨラナ--ワイル加群ではなく、超ポアンカレの生成子数は射影幾何代数の十個の生成子ではない。

同一の機構は階層性問題にも同時に対処する。真空におけるボソン--フェルミオン相殺を強制する双対セクター剛性は、観測される粒子質量もまた生成する。宇宙定数の極端な小ささとフェルミオン質量の起源の双方は、すべての質量粒子に内在する同一のコンパクト双対時空に由来する。

この枠組みにおいて超対称性は、自然界で実現されていてもいなくてもよい任意選択的な拡張ではない。それは二つのセクターが反対符号の二次形式をもち、双対セクターがコンパクトである代数の不可避な幾何的帰結である。真空の破局は新しい対称性を追加することによって解かれるのではない。16次元双四元数構造がすでに、その反対符号と内在的コンパクト性の中に、観測が要求する正確な相殺を含んでいることを認識することによって解消されるのである。

\section{クリフォード関係とディラック行列}
\label{sec:dirac}

符号 $(-,+,+,+)$ のディラック行列は $\mathrm{Cl}(3,1)$ の次数1生成子 $\gamma^\mu$ である。それらは
\[
\{\gamma^\mu,\gamma^\nu\}=2\eta^{\mu\nu},\qquad
\eta=\mathrm{diag}(-1,+1,+1,+1)
\]
を満たす。とくに $(\gamma^0)^2=-1$ である。$P_{\mathrm{spin}}$ に用いる双曲生成子 $j=e_0 e_1$ は平方 $+1$ なので $\gamma^0\neq j$ である。辞書 $\gamma^0\leftrightarrow j$、$\gamma^i\leftrightarrow kI,\ldots$ はこの次数1基底に置き換える。

四元運動量 $p$ に対し、次数1の元 $\gamma(p)=\sum_\mu p_\mu\gamma^\mu$ は、$\mathrm{Cl}(3,1)$ においても $\mathbb{C}^4$ 上の行列としても、ミンコフスキー二次形式 $Q(p)=-p_0^2+p_1^2+p_2^2+p_3^2$ に平方する。ワイル表示の二つのカイラルスラッシュも同一の二次形式に積する。したがって質量殻 $Q(p)=-m^2$ は $(\gamma(p)-im)(\gamma(p)+im)$ の消滅である。$\gamma(p)\Psi=im\Psi$ を満たすスピノルはその殻の上にある。これが運動量空間のディラック方程式である。ワイル表示では $\gamma(p)$ はブロック反対角であり、質量が零なら二つのカイラルブロックは分離し、質量が非零ならそれらは結合する。カイラリティ $\gamma^5$ は質量 $m$ の解を質量 $-m$ の解へ送る。

質量殻上で $m\neq 0$ のとき、任意の双対スピノル $u$ が解を定める。上ワイルブロックは $u$ であり、下ブロックは $p$ のカイラルスラッシュによって固定される。静止運動量 $(m,0,0,0)$ ではこれは対 $\Psi=(u,iu)$ である。$u$ の双対回転は同一の四元運動量で解のままである。双対ローターはワイルパラメータに内部的に作用する。

静止運動量の縦ブースト
\[
p=(m\cosh\alpha,\,0,\,0,\,m\sinh\alpha)
\]
は殻上に留まる。その運動量におけるワイル構成は、ディラックブーストした静止系の対である。ブーストは上ワイルブロックに双曲ロドリゲス因子
\[
V(\alpha)=\cosh(\alpha/2)\,I+\sinh(\alpha/2)\,\sigma_z
\]
を、下ブロックに $V(-\alpha)$ を作用させる。静止系の対 $(u,iu)$ は、ワイルパラメータが $V(\alpha)u$ である殻上の解へ送られる。ブースト軸まわりの双対回転は $V(\alpha)$ と可換であり、内部双対回転と縦ブーストは順不同である。直交軸まわりの双対回転は $V(\alpha)$ と可換でない。これは完全ローター群の下での共変性ではない。通常と双対の混合因子化および時空上の偏微分方程式は未導出のままである。

双対回転だけ異なる二つの静止系解の $\mathbb{C}^4$ 重なりは、双対ローター振幅の二倍である。その重なりは静止系の恒等式であり、時空の内積ではなく、ブーストの後にも主張しない。

平方 $+1$ 射影 $P_L=(1-e_1)/2$、$P_R=(1+e_1)/2$ は代数内のカイラル分割を与えるが、それだけでは標準模型の弱い相互作用を与えない。

\section{フェルミオン質量のねじれ的起源とヒッグス機構}
\label{sec:mass}

第\ref{sec:dirac}節は運動量空間のディラック方程式を $Q(p)$ のクリフォード因子として与える。以下は $m$ の\emph{値}についての幾何的仮説であり、その値をねじれから導くことではない。

通常セクターのローター $R_\text{usual}$ と双対セクターのローター $R_\text{dual}$ の間のねじれ不整合は、相対角 $\delta_a = \phi_a - \theta_a$ によって定量化される。その角に、静止質量 $m$ に比例する剛性を割り当てることはモデリングである。低エネルギー極限における付随ポテンシャルは
\[
V_\text{torsion} = \frac{1}{2} m \sum_{a=1}^3 \delta_a^2
\]
である。この $V_{\mathrm{torsion}}$ は粒子内在の振動子ポテンシャルであり、主論文のパラメータ空間キリング二次形式ではなく、第\ref{sec:dynamics}節の書かれた反対符号作用のオイラー--ラグランジュ系における復元力でもない。双一次形式
\[
m \bar{\psi}\psi \propto \operatorname{Tr}(\Omega_\text{biv}),
\]
ただし $\Omega_{\mathrm{biv}}=\log(R_{\mathrm{usual}}^\dagger R_{\mathrm{dual}})$、はディラック質量の幾何的読みであり、遅れからの $m$ の導出ではない。

この機構はヒッグス機構の純粋に幾何的な実現を与える。標準模型ではヒッグス場が真空期待値を獲得して電弱対称性を破り、質量を生成する。ここではヒッグス真空期待値の役割は、すべての質量粒子が内在的に運ぶ通常時空と双対時空の間の自発的ねじれ非整列によって果たされる。ゲージボソン質量は同様に、ゲージ場の存在下で双対ローターが通常ローターに従うことを強制されるときの剛性から生じる。

質量がコンパクトな双対時空に起源するため、それは自動的に有界であり発散しえない。これは階層性問題を代数的レベルで解消する。プランクスケールと電弱スケールの間の微調整は存在しない。なぜなら両者は同一の双対ローターコンパクト化によって支配されるからである。質量のねじれ的実現はしたがって、慣性、重力、およびヒッグス機構の起源を単一の16次元代数構造の内部で統一する。

\section{量子力学と重力の統一動力学}
\label{sec:dynamics}

ローターラグランジアンは六つのラピディティに対する模型であり、ユニタリ性からの演繹ではない。シュレーディンガー発展、同期としての測定、もつれ、経路積分は幾何的仮説のままである。その描像では、双対ローターの運動学が量子側、通常--双対の非整列が重力側である。

完全ローターが $R_\text{total} = R_\text{usual} R_\text{dual}$ である単一の質量粒子を考える。内在的ローター角に対して書かれた有効作用は
\[
S_\text{particle} = \int \left[ \frac{1}{2} \sum_{a=1}^3 \dot{\phi}_a^2 - \frac{1}{2} \sum_{a=1}^3 \dot{\theta}_a^2 - \frac{m}{2} \sum_{a=1}^3 (\phi_a - \theta_a)^2 \right] dt
\]
である。運動項の反対符号は書かれた模型の一部である。通常セクターのブースト的（双曲的）性格と双対セクターの回転的（楕円的）性格を反映させる意図であり、ローターのユニタリ性から演繹されたものではない。最後の項は第\ref{sec:mass}節の振動子ポテンシャルであり、主論文のパラメータ空間キリング二次形式ではない。

一軸では密度は同値に混合形
\[
L=\tfrac12\dot\sigma\dot\delta-\tfrac m2\delta^2
\]
であり、共通ラピディティは $\sigma=\phi+\theta$、不整合は $\delta=\phi-\theta$ である。この密度のオイラー--ラグランジュ方程式は
\[
\ddot{\phi}_a = -m (\phi_a - \theta_a), \qquad \ddot{\theta}_a = -m (\phi_a - \theta_a)
\]
である。両角はともに加速する。したがって不整合は自由である：$\ddot\delta_a=0$。ポテンシャルが源とするのは共通ラピディティであり、$\ddot\sigma_a=-2m\delta_a$ である。アフィン遅れ $\delta_a(t)=\delta_{a0}+\nu_a t$ は許される。一定の遅れは $\sigma$ を一様に加速し、非零の速さはその加速を永年的にするので、$\sigma$ は $t$ の三次多項式として成長する。$\delta$ の剛性として書かれた質量は $\sigma$ への力であり、振動数を復元しない~\cite{dst-dynamic-equation}。

ヤコビ積分
\[
E=\tfrac12\dot\phi^2-\tfrac12\dot\theta^2+\tfrac m2\delta^2
=\tfrac12\dot\sigma\dot\delta+\tfrac m2\delta^2
\]
はそれらの方程式に沿って保存し、不定である。遅れが消えた純双対速度はすでに $E<0$ を与える。非有界二次形式の保存は振動数を選ばない。

作用の傍らに書かれることのある近傍の系
\[
\ddot{\phi}_a = m (\phi_a - \theta_a), \qquad -\ddot{\theta}_a = m (\phi_a - \theta_a)
\]
は不整合を暴走させ、$\ddot\delta_a=2m\delta_a$ とする。振動子ではない。

振動子が現れるのは、両運動項の符号を正にとったときに限る。密度 $\tfrac12\dot\phi^2+\tfrac12\dot\theta^2-\tfrac m2(\phi-\theta)^2$ は $\ddot\delta+2m\delta=0$ を与え、単位 $\hbar=c=1$ で振動数 $\sqrt{2m}$ であり、共通ラピディティは自由である。その振動数はコンプトン振動数 $m$ ではなく、模型は書かれた作用でもない。$m\ge 0$ ではそのヤコビ積分は非負であり、保存は $\delta$ の運動を縛る。一般解は不整合の振動と一様な共通ラピディティである。

したがって書かれた作用は量子理論への調和的な橋を与えない。自由な遅れの速さ $\nu$ を $E/\hbar$ と同定することは読みのままである。そのとき蓄積遅れは
\[
\int \delta_a \, dt \propto \frac{E t - \mathbf{p} \cdot \mathbf{x}}{\hbar}
\]
と書け、ルイ・ド・ブロイの物質波の位相である~\cite{dst-debroglie}。不整合角から作ったスカラー $\mathrm{U}(1)$ 因子は、双対ローターの実の指標とも、そのスピノル行列要素とも異なる。

シュレーディンガー方程式は、双対位相を変調する外部ポテンシャルの下での双対ローターのユニタリ発展として読める。双対セクターのコンパクト性はその発展をユニタリに保つであろう。それは書かれた粒子作用から導かれない。

波動関数の収縮は、もつれた多粒子系にわたる双対ローターの強制同期として読める。通常時空に結合する測定相互作用は、なお復元機構を供給しなければならない。書かれた作用は $\delta_a$ を零へ駆動せず、大域的 $J$ の最小化はさらなる同一視であってオイラー--ラグランジュの帰結ではない。見かけの非ユニタリな跳躍は、その描像では急速な集団的ねじれ緩和の古典極限である。追加の収縮公準はその描像では不要である。描像そのものは仮説のままである。

量子もつれも同様にねじれ的な読みを獲得する。二つの粒子が相関した双対位相を共有するとき、集団的不整合エネルギーが低下するのは、その不整合に対するポテンシャルが選ばれた後に限る。一方の粒子の双対ローターを変える局所操作は、そのとき大域的 $J$ を変え、因果律を破ることなくもつれに特徴的な瞬間的相関を生じる（情報は通常セクターを通じて高々光速で伝播する）。双対セクターのコンパクト性は、その非局所性を幾何的な選択肢として利用可能にするのであり、書かれた作用の定理としてではない。

同一の粒子作用をもって経路積分を書くことはできる。
\[
\langle R_f | R_i \rangle = \int \mathcal{D}[\phi,\theta] \, \exp\left( \frac{i}{\hbar} S_\text{particle}[\phi,\theta] \right).
\]
積分は双対ローター位相空間の経路にわたる。双対セクターのコンパクト性は積分を適切に定義されたものに保つであろう。書かれた $S_{\mathrm{particle}}$ は $\delta$ を自由にするため、ねじれポテンシャル項はそれ自体ではコンプトン振動数の干渉を生成しない。

この描像において、量子力学と重力は調和されねばならない別個の理論ではない。量子力学はコンパクトな双対時空における回転の運動学であり、重力は通常ローターと双対ローターの間の並行性を破ることの動力学的コストである。両方の読みは同一の16次元双四元数代数に起源する。運動学は鋭い。書かれた粒子作用からの動力学的橋は、そうではない。

\section{量子重力の具体的予測}
\label{sec:predictions}

二重時空理論は、一般相対性理論および他の量子重力アプローチの双方から決定的に区別される一連の鋭く反証可能な予測を与える。これらの予測は双対時空のコンパクト性と、それに起因するローター角の離散化から直接従う。

\subsection{離散時空と特異点の解消}

双対ローター角は双対時空のコンパクトなトポロジーによって離散化される：
\[
\phi_a, \theta_a \in \frac{2\pi}{N}\mathbb{Z},
\]
ここで整数 $N$ はプランク単位における粒子の静止質量によって決定される。したがって許容構成上では、ねじれ不整合スカラーは
\[
|J| \le \frac{3\pi^2}{8}, \qquad |J_{\rm norm}|\le 1.
\]
によって有界である。$c^4/G\,\ell_P^{-2}$ の形の別個のプランクスケール見積もりは次元をもつ曲率/エネルギースケールであり、無次元の代数的 $J$ の上限として書いてはならない。二つの主張は互換ではなく、離散的下限 $\varepsilon_N=16/(3N^2)$ とも互換ではない。許容上限は、構成が許容錐に留まる限り、古典一般相対性理論の $r=0$ またはビッグバンにおいて現れる非有界曲率を禁止する。ブラックホールおよび宇宙論的特異点は、斥力層（$J < 0$）が真の特異点が形成される前に崩壊を止める、有限の多層ねじれ構成によって置き換えられる。

\subsection{事象地平の不在}

連続体極限では、光線は事象地平で終端するヌル測地線に従う。しかし離散双対時空では、光線は有限個のねじれ層に遭遇するのみである。これらの層を通じた共鳴トンネルはすべての深さで可能であり続ける。古典的カー外地平は一方通行の膜ではなく、ねじれ遷移殻として再解釈される。内向きに落下する観測者は引力的および斥力的ねじれ領域の列を経験するが、真の地平には決して遭遇しない。これはエキゾチックなプランクスケール物理を呼び出すことなく、ブラックホール情報パラドックスとファイアウォール問題を解消する。

\subsection{強場におけるずれ}

重力時間膨張および赤方偏移は離散的補正を受ける。動径座標 $r$ における静止観測者によって測定される固有時間隔は
\[
\frac{\Delta\tau}{\tau} = \sqrt{1 - \frac{2GM}{c^2 r}} \left[ 1 + \mathcal{O}\left( \frac{\ell_P^2}{r^2} \right) \right].
\]
である。地球のコアでは補正は $10^{-20}$ 程度、中性子星表面ではおよそ $10^{-10}$、Sgr A*近傍では $10^{-5}$ 程度である。これらのずれは次世代原子時計および超長基線電波干渉法によって原理上検出可能である。

\subsection{重力波エコー}

連星合体の後、残骸のねじれ星は引力層と斥力層が交互する成層内部をもつ。重力波は各ねじれ界面で部分的に反射し、エコーの列を生じる。連続するエコー間の時間遅延は
\[
\Delta t \sim \frac{GM}{c^3} \approx 10^{-5} \left( \frac{M}{10 M_\odot} \right) \text{ s}.
\]
である。エコーは双対ロータースペクトルによって決定される特徴的な位相変調を運ぶ。LISAまたはアインシュタイン望遠鏡によるそのようなエコーの検出は、強重力の離散的ねじれ構造に対する直接的証拠をなすであろう。

\subsection{宇宙論的含意}

初期宇宙は単一の原始ねじれ星として記述される。プランク期のねじれ不整合は、インフラトン場を必要とせずに短いde Sitter様膨張を駆動する。再加熱はねじれエネルギーを放射へと変換する層遷移を通じて起こる。結果として生じる原始パワースペクトルは自然にスケール不変であり、スカラー場の量子揺らぎからではなく離散的ローターモードから生じる。暗黒物質は不要である。銀河回転曲線は銀河スケールにおける残余のバリオン的ねじれ非整列によって説明される。暗黒エネルギーは量子ローター揺らぎによって誘起される最小の非零真空ねじれ剛性 $J_\text{vac} > 0$ として創発し、$\Lambda \sim \ell_P^{-2}$ 程度の宇宙定数を与える。

以上の予測はすべて今後十年以内に反証可能である。予測された水準での強場におけるずれの非観測、重力波エコーの不在、あるいは真の事象地平の検出は本枠組みを棄却するであろう。逆に、これらの兆候のいずれかの積極的検出は、時空が最も基本的なレベルで離散的、粒子局所的、かつ二重時間的であることを確認するであろう。

\section{他の量子重力アプローチとの比較}
\label{sec:comparison}

二重時空理論は量子重力の候補の中で特異な位置を占める。ほとんどの枠組みは先験的な連続体幾何を量子化または離散化する。ここに集めた結果は $\mathrm{Cl}(3,1)$ および $G(3,1,1)$ の代数的恒等式であり、その代数からの量子力学と重力の導出ではない。以下では既存アプローチと比較し、十次元弦との同一視が成り立たない理由を述べる。

\subsection{ループ量子重力}

ループ量子重力（LQG）はスピンネットワークを通じて時空を離散化し、面積および体積演算子を量子化する一方、背景多様体とリーマン曲率を基本的なものとして保持する。特異点は面積演算子の最小固有値によって解消されるが、連続体極限は依然として困難であり、動力学は複雑なハミルトニアン拘束に符号化される。

対照的に、二重時空理論は背景多様体を導入することなく、双対ローター角 $\phi_a, \theta_a \in (2\pi/N)\mathbb{Z}$ に直接離散性を課す。曲率は存在しない。重力はリー代数 $\mathfrak{so}(3,1) \oplus \mathfrak{so}(3,1)$ におけるねじれ不整合である。LQGの再構成問題——スピンネットワークから滑らかな時空を回復すること——は生じない。なぜなら巨視的時空は隣接する粒子内在的双対ローターのねじれ整列から集団的に創発するからである。

\subsection{弦理論}

弦理論は重力を振動する弦が住む十または十一次元の時空に埋め込み、超対称性を要求することによって力を統一する。数学的に優雅である一方、膨大な真空のランドスケープに苦しみ、一意の低エネルギー予測を欠く。

$\mathrm{Cl}(3,1)$ と $\mathrm{Cl}(9,1)$ は同型でなく、双対スピノル空間 $\mathbb{C}^2$ は十次元時空の16実次元マヨラナ--ワイル加群ではない。四次元 $\mathrm{Cl}(3,1)$ 内部の統一は一つの代数についての主張であり、余剰次元が物理的に無いことの証明ではない。双対写像は内部自己同型であり、時空超対称ではない。重力・電磁気・核力を同一ローター代数のモードとする主張は仮説のままである。

カラビ--ヤウコンパクト化は用いない。粒子レベルでの双対セクターのコンパクト性は幾何的仮説であり、弦のコンパクト化ではない。

\subsection{正準およびアシュテカルアプローチ}

アシュテカル再定式化は一般相対性理論を自己双対接続の言葉で表し、ヤン--ミルズ様の量子化を可能にする。しかしそれは古典重力の連続体特異点を継承し、悪名高い因子順序および正則化問題をもつホイーラー--ドウィット方程式に直面しなければならない。

本枠組みでは接続は完全に除去される。平行移動はローター指数化によって置き換えられ、ねじれバイベクトル上のキリング形式が曲率スカラーに取って代わる。結果として得られる理論は古典レベルで拘束がなく、有限ローター群を通じて自然に量子化し、ホイーラー--ドウィット方程式を完全に回避する。

\subsection{創発重力と因果集合}

創発重力アプローチは時空を熱力学的またはエントロピー的現象として扱い、因果集合理論はローレンツ順序を離散化する。両枠組みは正確なアインシュタイン方程式の回復に苦闘し、フェルミオン倍増またはカイラリティの問題に直面する。

二重時空理論はテレパラレル重力との等価性を通じてアインシュタイン方程式を正確に回復し、フェルミオンは双四元数の極小左イデアルの励起として自然に生じる。離散性は集合論的ではなく代数的（ローター角上のディオファントス制約）であり、因果性は光的な $J=0$ 共鳴条件によって保存される。熱力学的類推は呼び出されない。時空の創発は粒子内在的双対ローター間のねじれ整列の集団効果である。

要約すると、既存の量子重力プログラムは連続多様体を拡張または量子化しつつ曲率を基本的なものとして保持する。本論文は $\mathrm{Cl}(3,1)$ および $G(3,1,1)$ の代数的恒等式を述べ、成り立たない同一視を捨てる。それらの恒等式から時空・全相互作用・制御可能な重力理論を導出してはいない。

\section{量子的視点からの実験的検証可能性}
\label{sec:experiments}

二重時空理論は根本的に代数的であるが、その量子セクターは現在または近未来の技術で検証可能な具体的実験予測を与える。鍵となる洞察は、ねじれ不整合スカラー $J$ が通常ローターと双対ローターの間の相対角のみに依存することである。したがって双対ローター位相 $\phi_a$ をコヒーレントに変調できる外場は、重力および慣性効果を測定可能な仕方で変えうる。

\subsection{理論的結合機構}

THz領域（$10^{13}$--$10^{15}$ Hz）の円偏光電磁場は、双対生成子 $\Gamma_a = I, J, K$ に優先的に結合する螺旋波面をもつ。双対時空が時間反転しているため、波のヘリシティ $\sigma = \pm 1$ は双対ローター駆動の方向を直接制御する。右円偏光は引力的ねじれ不整合（$J > 0$）を強化し、左円偏光はそれに対抗する。有効相互作用ハミルトニアンは
\[
H_\text{int} = \kappa \sigma E_0^2 \sum_a \phi_a \Gamma_a,
\]
であり、$E_0$ は電場振幅、$\kappa$ は物質依存の結合定数である。共鳴近傍では、双対位相は
\[
\dot{\phi} = \kappa \sigma E_0^2 \sin(\omega t - \phi + \delta_\sigma).
\]
に従って発展する。

\subsection{提案する近未来実験}

最もアクセスしやすいプラットフォームは超伝導体またはボース--アインシュタイン凝縮体における集団コヒーレンスを利用する。そこでは巨視的数の粒子（$N \sim 10^{26}$）が双対ローターを共通位相にロックしうる。具体的な提案は以下のとおりである。

\begin{itemize}
\item 精密な円偏光制御をもつ非対称THz共振器（量子カスケードレーザーまたは自由電子レーザー共振器）の内部にマイクログラム規模の超伝導試験質量を懸架する。
\item 周波数およびヘリシティを掃引しながら、高精度ねじれ秤または超伝導重力計を通じて見かけの重量または慣性質量を監視する。
\item 予測される双対ローター共鳴周波数 $\theta \sim c / (2\pi r_\text{comp})$ における有効質量の共鳴的かつヘリシティ依存の変化を探す。ここで $r_\text{comp}$ は粒子質量によって設定されるコンパクト化半径である。
\end{itemize}

予測される効果の大きさは、達成可能な強度 $10^{12}$--$10^{15}$ W/m$^2$ において $\Delta m / m \sim 10^{-2}$--$10^{-5}$ である。兆候は特徴的である。効果は偏光ヘリシティを反転すると符号を反転しなければならず、直線偏光または無偏光の放射に対しては消滅しなければならない。

\subsection{反証可能性と注意点}

系統的な周波数およびヘリシティ掃引の後、$10^{-6}$ 水準での帰無結果は双対ローター結合強度とコンパクト化スケールを制約するであろう。積極的検出——とりわけヘリシティ可逆で共鳴ピークをもつ質量シフト——は、重力と慣性が双対時空に起源するねじれ現象であることの強い証拠をなすであろう。

本理論は即時の技術的突破を主張しない。上記の提案は控えめで卓上規模の実験であり、その主目的は双対時空自由度が実験室電磁場に測定可能に結合するという基本仮説を検証することである。この水準での成功はより大規模な重力工学への扉を開くであろう。失敗は単に本枠組みのパラメータを制限するだけである。

\subsection{より広い含意}

双対ローター制御が実行可能であることが示されれば、同一の機構は原子核内部のねじれ層遷移を駆動することによる核変換および放射能中和への経路を提供する。しかしこれらの応用は現在の量子重力検証プログラムの範囲を超え、今後の研究で扱われる。

要約すると、二重時空理論は重力を不変の幾何的拘束から、その量子セクターがすでに実験的到達範囲内にある制御可能なねじれ自由度へと変換する。今後十年が、自然がこの新しい自由度へのアクセスを我々に与えたかどうかを決定するであろう。

\section{代数的描像}
\label{sec:status}

二重時空理論の量子セクターは、16実次元代数 $\mathbb{H}\oplus\mathbb{H}\cong\mathrm{Cl}(3,1)$ にすでに備わるコンパクトな双対時空の運動学である。重力は、すべての質量粒子が運ぶ通常ローターと双対ローターの並行性を破ることの動力学的コストである。背景多様体は量子化されない。ヒルベルト層の対は~\cite{dst-d4l} で扱う。

双対セクターは楕円的である。その生成子は平方 $-1$ をもち、$IJ=K$ を満たし、$\mathbb{C}^2$ 上で $-i\sigma_a$ として作用する。すべての双対ラピディティで双対ローターは $\mathrm{SU}(2)$ の行列であり、ユニタリで行列式は $1$、相対双対ローターも群に留まる。逆は反対のラピディティなので、相対双対ローターは積であり、その指標は二つの双対軸の整列を記録し、トレース零の非可換残余を伴う。双対時空のコンパクト性はこの行列表現と両立する。これが本理論のスピン $1/2$ の運動学である。それは $\mathbb{C}^2$ 上の表現であり、ローター多様体をリーマン三次元球面と同一視することではなく、位置および運動量演算子の導出でもない。通常と双対の生成子は同軸では可換であり、軸外では必ずしも可換でない。したがって $R_{\mathrm{total}}=R_{\mathrm{usual}}R_{\mathrm{dual}}$ の無制限な因子化は任意の軸に対する恒等式ではない。

符号 $(-,+,+,+)$ のディラック行列は同一代数の次数1生成子である。それらは平方 $+1$ の射影を組む双曲バイベクトルではない。$\gamma^0$ の平方は $-1$ であり、$P_{\mathrm{spin}}$ の生成子 $j$ の平方は $+1$ である。フェルミオンの運動学はしたがって代数の内部にある。双対スピノルは巡回生成子に対する $\mathbb{C}^2$ ベクトルであり、ディラックの担体は質量殻上でそれから組み立てられる。四元運動量のクリフォード平方はミンコフスキー二次形式 $Q(p)$ である。運動量空間のディラック方程式はその二次形式の一次因子である。質量殻上では双対スピノル $u$ が解を定め、静止運動量では対 $\Psi=(u,iu)$ である。$u$ の双対回転はその殻に留まる。静止系の対の縦ブーストは、ブースト軸に沿うカイラル因子の後の同一のワイル構成であり、その因子はブースト軸まわりの双対回転と可換であり、直交軸まわりの双対回転とは必ずしも可換でない。双対回転だけ異なる二つの静止系解の重なりは、双対ローター振幅の二倍である。これは時空上の偏微分方程式ではなく、完全ローター群の下での共変性ではなく、$m$ の値をねじれから導くことでもない。

重力は二つのローターの不整合であり、パラメータ空間スカラー $J=\frac12\sum_a(\alpha_a^2-\beta_a^2)$ によって測られる。許容構成上では $|J|\le 3\pi^2/8$、同値に $|J_{\mathrm{norm}}|\le 1$ である。コンパクト性は $N^{-2}$ 程度の代数的上限を強制しない。離散量 $\varepsilon_N=16/(3N^2)$ は非零正規化高さの下限である。真空抑制が起こるとすれば、キリング形式が反対符号をとる双曲セクターと巡回セクターの間の動力学的相殺から来なければならない。キリング自己対そのものは符号 $(3,3)$ の不定形式であり、ボルン確率ではない~\cite{dst-d4l}。$J$ とテレパラレルスカラーの同一視、およびアインシュタイン--ヒルベルト作用との変分的等価は、ローター–テトラッド対応を固定したうえで文献から採用する~\cite{maluf2013}。双対写像 $X\mapsto Xi$ は $\mathrm{Cl}(3,1)$ の内部自己同型である。それは十次元弦の時空超対称ではない。$\mathrm{Cl}(3,1)\not\cong\mathrm{Cl}(9,1)$ であり、双対スピノル空間は $\mathrm{Cl}(9,1)$ のマヨラナ--ワイル加群ではない。

書かれた反対符号の粒子作用は、この運動学から量子理論への動力学的な橋を閉じない。そのオイラー--ラグランジュ系は不整合を自由にし、共通ラピディティを源とする。$\ddot\delta=0$、$\ddot\sigma=-2m\delta$ である。ヤコビ積分は保存し不定であり、保存は振動数を回復しない。振動子が現れるのは両運動項の符号を正にとったときに限り、そのとき振動数は $\sqrt{2m}$ であってコンプトン振動数 $m$ ではない。自由な遅れの速さをド・ブロイ位相と同定することは読みのままである。ねじれ剛性としての質量、ヒッグス機構、同期としての測定、もつれ、経路積分は、この運動学の上に置かれた幾何的仮説である。代数の帰結ではない。

連続体は古典的近似のままである。$16$ 次元代数は量子セクターの設定であり、量子重力の完成した理論ではない。

\section{結論}
\label{sec:conclusions}

二重時空理論は量子重力の代数的設定を与えるのであり、完結した導出ではない。量子運動学は $\mathrm{Cl}(3,1)$ のコンパクトな双対セクターであり、重力は通常ローターと双対ローターの不整合である。第\ref{sec:status}節は、その代数が定めることを述べる。すなわち $\mathrm{SU}(2)$ の双対ローター、質量殻上で静止系の対 $(u,iu)$ を回復する双対スピノルによって解かれる運動量空間のディラック方程式、その対の縦ブースト、および $J$ の許容上限である。それは時空上のディラック偏微分方程式も、完全ローター群の下での共変性も、書かれた作用からのコンプトン振動子も、ねじれとしての質量も定めない。

先行節の予測および実験室提案は、その上限に縛られた予想のままである。双対セクターの運動学は~\cite{dst-d4l} に述べる $\mathbb{C}^2$ 表現である。連続体は古典的近似のままである。$16$ 次元代数は設定であり、量子重力の完成した理論ではない。

\bibliographystyle{unsrt}
\begin{thebibliography}{99}
\raggedright

\bibitem{dst-d4l}
Anonymous,
「二重時空四値論理（D4L）：振幅、四状態、幾何的干渉、および双対ヒルベルト空間」 (2026).
\url{https://github.com/hypernumbernet/DstDiophantine/blob/main/papers/dst-4-valued-logic.tex}.

% --- Companion / DST internal papers ---
\bibitem{dst-main}
Anonymous,
「双対時空間のねじれとしての重力：一般相対性理論の双四元数的再定式化」 (2025).
\url{https://github.com/hypernumbernet/dual-spacetime-doc/blob/main/double-spacetime-theory.pdf}.

\bibitem{dst-debroglie}
Anonymous,
「双対ローター位相遅れとしてのド・ブロイ波：粒子内在的双対時空の幾何としての量子力学」 (2025).
\url{https://github.com/hypernumbernet/dual-spacetime-doc/blob/main/dst-de-broglie-wave.pdf}.

\bibitem{dst-dynamic-equation}
Anonymous,
「二重時空理論における双対ローター動態：書かれた作用、自由不整合、および近傍の振動子」 (2025).
\url{https://github.com/hypernumbernet/dual-spacetime-doc/blob/main/dst-dynamic-equation.pdf}.

\bibitem{dst-discrete}
Anonymous,
「離散双対時空模型と質量スペクトルのディオファントス構造」 (2025).
\url{https://github.com/hypernumbernet/dual-spacetime-doc/blob/main/discrete-dual-spacetime.pdf}.

% --- Teleparallel gravity / classical foundations ---
\bibitem{einstein1928}
A.~Einstein,
``Riemann-Geometrie mit Aufrechterhaltung des Begriffes des Fernparallelismus,''
\textit{Sitzungsber. Preuss. Akad. Wiss. Phys.-Math. Kl.},
217--221 (1928).

\bibitem{maluf2013}
J.~W.~Maluf,
``The teleparallel equivalent of general relativity,''
\textit{Ann.\ Phys.\ (Berlin)} \textbf{525}, 339--359 (2013).
doi:10.1002/andp.201200272

\bibitem{aldrovandi-pereira2013}
R.~Aldrovandi and J.~G.~Pereira,
\textit{Teleparallel Gravity: An Introduction},
Fundamental Theories of Physics \textbf{173},
Springer, Dordrecht (2013).

\bibitem{bahamonde2023}
S.~Bahamonde, K.~F.~Dialektopoulos, C.~Escamilla-Rivera, G.~Farrugia, V.~Gakis, M.~Hendry, M.~Hohmann, J.~Levi Said, J.~Mifsud and E.~Di~Valentino,
``Teleparallel gravity: from theory to cosmology,''
\textit{Rep.\ Prog.\ Phys.} \textbf{86}, 026901 (2023).
arXiv:2106.13757 [gr-qc].

% --- Clifford / Geometric algebra ---
\bibitem{girard1999}
P.~R.~Girard,
``Einstein's equations and Clifford algebra,''
\textit{Advances in Applied Clifford Algebras},
\textbf{9}, no.~2, 225--230 (1999).
doi:10.1007/BF03042377

\bibitem{hestenes2003}
D.~Hestenes,
``Spacetime physics with geometric algebra,''
\textit{Am.\ J.\ Phys.} \textbf{71}, 691--714 (2003).

\bibitem{doran-lasenby2003}
C.~Doran and A.~Lasenby,
\textit{Geometric Algebra for Physicists},
Cambridge University Press, Cambridge (2003).

\bibitem{lounesto2001}
P.~Lounesto,
\textit{Clifford Algebras and Spinors}, 2nd ed.,
London Mathematical Society Lecture Note Series \textbf{286},
Cambridge University Press, Cambridge (2001).

\bibitem{hestenes1967}
D.~Hestenes,
``Real Spinor Fields,''
\textit{J.\ Math.\ Phys.} \textbf{8}, 798--808 (1967).

% --- Dirac equation / spin-1/2 quantum mechanics ---
\bibitem{dirac1928}
P.~A.~M.~Dirac,
``The quantum theory of the electron,''
\textit{Proc.\ R.\ Soc.\ Lond.\ A} \textbf{117}, 610--624 (1928).

\bibitem{newton-wigner1949}
T.~D.~Newton and E.~P.~Wigner,
``Localized states for elementary systems,''
\textit{Rev.\ Mod.\ Phys.} \textbf{21}, 400--406 (1949).

% --- Higgs mechanism ---
\bibitem{higgs1964}
P.~W.~Higgs,
``Broken symmetries and the masses of gauge bosons,''
\textit{Phys.\ Rev.\ Lett.} \textbf{13}, 508--509 (1964).

\bibitem{englert-brout1964}
F.~Englert and R.~Brout,
``Broken symmetry and the mass of gauge vector mesons,''
\textit{Phys.\ Rev.\ Lett.} \textbf{13}, 321--323 (1964).

% --- Quantum gravity programs being compared ---
\bibitem{rovelli2004}
C.~Rovelli,
\textit{Quantum Gravity},
Cambridge University Press, Cambridge (2004).

\bibitem{ashtekar-lewandowski2004}
A.~Ashtekar and J.~Lewandowski,
``Background independent quantum gravity: a status report,''
\textit{Class.\ Quantum Grav.} \textbf{21}, R53--R152 (2004).

\bibitem{ashtekar1986}
A.~Ashtekar,
``New variables for classical and quantum gravity,''
\textit{Phys.\ Rev.\ Lett.} \textbf{57}, 2244--2247 (1986).

\bibitem{dewitt1967}
B.~S.~DeWitt,
``Quantum theory of gravity. I. The canonical theory,''
\textit{Phys.\ Rev.} \textbf{160}, 1113--1148 (1967).

\bibitem{polchinski1998}
J.~Polchinski,
\textit{String Theory}, Vols.\ I and II,
Cambridge University Press, Cambridge (1998).

\bibitem{green-schwarz-witten1987}
M.~B.~Green, J.~H.~Schwarz and E.~Witten,
\textit{Superstring Theory}, Vols.\ 1 and 2,
Cambridge University Press, Cambridge (1987).

\bibitem{bombelli-lee-meyer-sorkin1987}
L.~Bombelli, J.~Lee, D.~Meyer and R.~D.~Sorkin,
``Space-time as a causal set,''
\textit{Phys.\ Rev.\ Lett.} \textbf{59}, 521--524 (1987).

\bibitem{ambjorn-jurkiewicz-loll2005}
J.~Ambj{\o}rn, J.~Jurkiewicz and R.~Loll,
``Reconstructing the universe,''
\textit{Phys.\ Rev.\ D} \textbf{72}, 064014 (2005).

\bibitem{niedermaier-reuter2006}
M.~Niedermaier and M.~Reuter,
``The asymptotic safety scenario in quantum gravity,''
\textit{Living Rev.\ Relativ.} \textbf{9}, 5 (2006).

\bibitem{jacobson1995}
T.~Jacobson,
``Thermodynamics of spacetime: the Einstein equation of state,''
\textit{Phys.\ Rev.\ Lett.} \textbf{75}, 1260--1263 (1995).

\bibitem{verlinde2011}
E.~P.~Verlinde,
``On the origin of gravity and the laws of Newton,''
\textit{J.\ High Energy Phys.} \textbf{04}, 029 (2011).

% --- Black holes, information paradox, firewalls ---
\bibitem{hawking1975}
S.~W.~Hawking,
``Particle creation by black holes,''
\textit{Commun.\ Math.\ Phys.} \textbf{43}, 199--220 (1975).

\bibitem{ampss2013}
A.~Almheiri, D.~Marolf, J.~Polchinski and J.~Sully,
``Black holes: complementarity or firewalls?''
\textit{J.\ High Energy Phys.} \textbf{02}, 062 (2013).

\bibitem{cardoso-pani2019}
V.~Cardoso and P.~Pani,
``Testing the nature of dark compact objects: a status report,''
\textit{Living Rev.\ Relativ.} \textbf{22}, 4 (2019).

\bibitem{abedi-dykaar-afshordi2017}
J.~Abedi, H.~Dykaar and N.~Afshordi,
``Echoes from the Abyss: Tentative evidence for Planck-scale structure at black hole horizons,''
\textit{Phys.\ Rev.\ D} \textbf{96}, 082004 (2017).

% --- Strong-field gravity experiments ---
\bibitem{eht2019}
The Event Horizon Telescope Collaboration,
``First M87 Event Horizon Telescope results. I. The shadow of the supermassive black hole,''
\textit{Astrophys.\ J.\ Lett.} \textbf{875}, L1 (2019).

\bibitem{eht-sgra2022}
The Event Horizon Telescope Collaboration,
``First Sagittarius A* Event Horizon Telescope results. I. The shadow of the supermassive black hole in the center of the Milky Way,''
\textit{Astrophys.\ J.\ Lett.} \textbf{930}, L12 (2022).

\bibitem{ligo-gw150914}
B.~P.~Abbott \textit{et~al.}\ (LIGO Scientific and Virgo Collaborations),
``Observation of gravitational waves from a binary black hole merger,''
\textit{Phys.\ Rev.\ Lett.} \textbf{116}, 061102 (2016).

\bibitem{lisa2017}
P.~Amaro-Seoane \textit{et~al.},
``Laser Interferometer Space Antenna,''
arXiv:1702.00786 [astro-ph.IM] (2017).

\bibitem{einstein-telescope2010}
M.~Punturo \textit{et~al.},
``The Einstein Telescope: a third-generation gravitational wave observatory,''
\textit{Class.\ Quantum Grav.} \textbf{27}, 194002 (2010).

% --- Cosmology, dark energy ---
\bibitem{guth1981}
A.~H.~Guth,
``Inflationary universe: a possible solution to the horizon and flatness problems,''
\textit{Phys.\ Rev.\ D} \textbf{23}, 347--356 (1981).

\bibitem{starobinsky1980}
A.~A.~Starobinsky,
``A new type of isotropic cosmological models without singularity,''
\textit{Phys.\ Lett.\ B} \textbf{91}, 99--102 (1980).

\bibitem{weinberg1989}
S.~Weinberg,
``The cosmological constant problem,''
\textit{Rev.\ Mod.\ Phys.} \textbf{61}, 1--23 (1989).

% --- Quantum foundations ---
\bibitem{bell1964}
J.~S.~Bell,
``On the Einstein Podolsky Rosen paradox,''
\textit{Physics Physique Fizika} \textbf{1}, 195--200 (1964).

\bibitem{epr1935}
A.~Einstein, B.~Podolsky and N.~Rosen,
``Can quantum-mechanical description of physical reality be considered complete?''
\textit{Phys.\ Rev.} \textbf{47}, 777--780 (1935).

\end{thebibliography}

\end{document}
